
\section{A reporting standard, and software}
\label{sec:standard}

\subsection{The minimum reportable form}

The apparatus of Sections~\ref{sec:surface} and~\ref{sec:inference} implies a
reporting standard, and it is short enough to state in full. A feature-value
claim is reportable when it carries:

\begin{enumerate}
\item \textbf{the design space}, as an enumerable declaration --- which
learners, encodings, targets, splits, decision times, quality mechanisms,
baselines, metrics and operating points are admissible, why anything
defensible was left out, and \textbf{the measure placed on those levels};
\item \textbf{the denominator}, as a count that a reader can reproduce from
the axis levels: declared cells, computed cells, and the cells any inference
actually rests on, which need not be the same set;
\item \textbf{the absolute increments} at every rung of the baseline ladder,
in the metric's own units, with intervals from a resampling scheme that
refits the model;
\item \textbf{the sensitivity decomposition}, with the total higher-order
share reported as $1 - \sum_i S_i$ rather than as any per-axis difference;
\item \textbf{the resolution region and $\rho$}, under a band whose family is
the set the label ranges over;
\item \textbf{the ratio $R$ last if at all}, with a Fieller set and the set's
kind, and never where the denominator is not resolvably positive.
\end{enumerate}

Nothing in that list is expensive once the surface exists. The cost is the
surface, and the surface is a factorial over declarations the analyst has
already made implicitly.

\subsection{\texttt{fieldvalue}}
\label{sec:tool}

The standard ships as \texttt{fieldvalue}, a package that computes the surface
from a data frame, a feature name and a dictionary of baselines. Its notable
property is a refusal: \texttt{Surface} raises rather than returning a scalar.
\texttt{Surface.report()} emits the minimum reportable form above;
\texttt{Surface.at(...)} returns one cell and requires the caller to name every
axis. Version \fvVersion\ adds the four objects of
Sections~\ref{sec:sobol}--\ref{sec:regretdef} as functions of a surface ---
\texttt{decompose}, \texttt{regions}, \texttt{robustness}, \texttt{regret} ---
each taking a tidy frame rather than this package's own object, so a reader can
apply them to a surface computed by some other pipeline.

The package carries \nTests\ tests. It is also an \emph{independent
implementation}: the analysis scripts and the package compute the same four
objects from the same definitions in different code, neither written against
the other, and \texttt{scripts/s11\_tool\_agreement.py} runs both on the same
input. On this paper's corpus they agree on \nAgreeQuantities\ of
\nAgreeQuantities\ quantities, with a largest difference of \agreeTolerance.

\subsection{Beyond process logs, and beyond this paper's claims}
\label{sec:worked}

The standard is not about event logs and not about configuration databases.
Supplement~\ref{app:worked} runs it on UCI Adult, where the expensive field is
\texttt{occupation} and the cheap ones come with the sampling frame: over
\exCells\ cells the reduction runs \exRmin\ to \exRmax; the baseline and the
metric carry the same share of the variance to three figures,
\exSbaselinePct\ and \exSmetricPct, against \exSpopulationPct\ for the
register population; and a one-number report misstates the sign for a mean
\exMisreportPct\ of the other cells. \texttt{examples/worked\_example.py}
produces all of it in \exampleLines\ statements against the package's public
interface.

\paragraph{How the surface is reported today}
\label{sec:practice}
Whether the axes this paper varies are varied in published work deserves a
systematic review, and we did not conduct one. A machine-assisted prevalence
pilot was run on a convenience frame and is \emph{not} reported here: with one
adjudicator, no independent human reliability study and a mechanical screen
whose measured sensitivity is \auditSens, it can generate a hypothesis about
practice and cannot test one. Its frame, sample, codes, adjudications and
evidence dossiers are in the archive. The paper's argument does not need it:
the case against one-number reporting is made by the decomposition and the
sign disagreement, on data, and not by a count of how often other authors
report one number.
